\documentclass[../main.tex]{subfiles}

\begin{document}

\chapter*{Preface}
\addcontentsline{toc}{chapter}{Preface}

\section*{Optimization: The Heartbeat of Machine Learning}

Every time a neural network learns, optimization is at work. Every weight update, every training step, every improvement in model performance---all are driven by optimization algorithms. Without optimization, machine learning would be nothing more than static equations on a page. Optimization breathes life into our models.

Yet optimization is often treated as a black box. Engineers call \texttt{optimizer.step()} without understanding why Adam outperforms vanilla SGD on their problem, or why their training suddenly diverges. They tune learning rates by trial and error, never grasping the theory that could guide their decisions. They fight overfitting with regularization tricks they don't fully understand.

This book changes that. Here, you'll develop a deep, practical understanding of optimization---the mathematical engine that powers all of machine learning.

\section*{What This Book Covers}

This book presents seven core chapters plus a hands-on capstone project on optimization for machine learning, building from theoretical foundations to practical algorithms:

\textbf{Chapter 1: Convex Optimization} establishes the mathematical foundation. You'll learn what makes a problem convex, why convexity guarantees we can find global optima, and how these concepts apply to machine learning. Even when deep learning problems aren't convex, understanding convexity helps you recognize when optimization will be easy and when it will be hard.

\textbf{Chapter 2: Stochastic Optimization Methods} introduces the workhorse algorithm of deep learning. You'll understand why using random minibatches works, how the noise in SGD can actually help escape local minima, and the theoretical guarantees that explain SGD's remarkable success.

\textbf{Chapter 3: Gradient Descent Variants} explores gradient descent and momentum methods. From basic gradient descent through momentum and Nesterov acceleration, you'll build intuition for how these algorithms navigate loss landscapes and when each variant excels.

\textbf{Chapter 4: Advanced Optimizers} covers the adaptive optimizers that dominate modern practice, along with learning rate schedules. You'll understand why Adam works so well, its failure modes, and the newer variants like AdamW that address its limitations. You'll also learn about learning rate warmup, decay strategies, and how to choose Adam versus SGD with momentum.

\textbf{Chapter 5: Generalization and Regularization} tackles the art of controlling model complexity. From L1 and L2 regularization to dropout and early stopping, you'll understand these techniques as optimization constraints and develop intuition for when each is appropriate.

\textbf{Chapter 6: Loss Functions and Optimization Objectives} examines how we define ``success'' for our models. Different tasks require different loss functions, and your choice profoundly affects both what the model learns and how optimization behaves. You'll understand the connections between loss functions, probability distributions, and training dynamics.

\textbf{Chapter 7: Convergence Analysis} brings theory and practice together. You'll learn to analyze whether your training is converging, diagnose common problems, and make principled decisions about hyperparameters based on convergence guarantees.

\section*{Who This Book Is For}

This book serves machine learning practitioners and researchers who want to move beyond treating optimization as a black box:

\textbf{ML engineers} who want to understand why their training behaves the way it does, and how to fix it when it doesn't.

\textbf{Researchers} who need to design custom optimizers or loss functions for novel architectures and tasks.

\textbf{Students} preparing for advanced coursework or industry positions who want solid foundations in optimization.

\textbf{Anyone} who has called \texttt{model.fit()} and wondered what's really happening inside.

\section*{Prerequisites}

This book assumes you have completed \textit{Book 1: The Math That Powers AI} or have equivalent background in:
\begin{itemize}
    \item \textbf{Linear algebra:} vectors, matrices, eigenvalues, and basic matrix operations
    \item \textbf{Calculus:} derivatives, gradients, the chain rule, and partial derivatives
    \item \textbf{Probability:} basic distributions, expectations, and variance
    \item \textbf{Convexity basics:} the definition of convex functions and sets
\end{itemize}

If you need a refresher, the Math Recap chapter provides a concise review of the essential concepts from Book 1 that we'll build upon.

\section*{Your Journey Through The AI Engineer's Library}

With this book, you complete the mathematical foundations of the series:

\textbf{Book 1: The Math That Powers AI} provided the mathematical toolkit---linear algebra, calculus, probability, and basic optimization concepts.

\textbf{Book 2: Optimization for AI} (this volume) deepens your optimization expertise, giving you mastery over the algorithms that train every modern AI system.

From here, you're ready to advance to the architecture and application books:

\textbf{Book 3: The Generative Revolution} shows how these optimization techniques apply to training neural networks, from simple perceptrons to modern transformers and diffusion models.

\textbf{Book 4: Inside the Black Box} explores interpretability, alignment, and understanding what neural networks really learn.

\textbf{Book 5: The Age of AI Agents} addresses autonomous AI systems that can reason, plan, and act in the world.

The journey from mathematical foundations to production systems is long but rewarding. Let's continue.

\vspace{1cm}
\noindent\textit{Dr.\ Vijay Raghavan}\\
\textit{2026}

\end{document}
